\chapter{项目进度与 4--6 周监督计划}
\label{chap:progress}

截至 2026-04-28，本报告对应的仓库状态是一个实验材料充足但工作区较脏的研究工作区。已经完成的部分包括：实验数据归档目录已创建，manifest 与 checksum 已生成并通过校验；报告工作区已初始化并渲染；Protocol-A、Protocol-C 和听测统计均有本地可读产物；远端 \texttt{origin/master} 的 P0 报告与脚本已被抽取到归档目录；主结果表已经升级为带 bootstrap confidence intervals 与 paired significance 的论文表。尚未完成的部分包括：P0 原始 outputs 未找回；Protocol-A 多个口径尚需冻结；归档目录和报告目录仍未 stage；部分实验结果文件处于 ignored 或 untracked 状态；NeurIPS 补实验的统一 logger schema 尚未实现。

\begin{table}[H]
	\centering
	\caption{当前项目进度与证据状态。}
	\label{tab:project-progress}
	\small
	\resizebox{\textwidth}{!}{%
	\begin{tabular}{lll}
		\toprule
			事项 & 当前状态 & 下一步 \\
		\midrule
		实验归档目录 & 已生成，checksum 通过 & 决定是否纳入版本控制。 \\
		报告工作区 & 已生成并写入学术论文式正文 & 继续审阅并按需要修订。 \\
		Protocol-A 主证据 & 204 条 audio-grouped 结果存在 & 冻结为唯一主文口径。 \\
		主结果统计表 & 已补 mean [95\% CI] 与 Holm-corrected paired tests & 与 logger run id 绑定。 \\
		Protocol-C 诊断 & 211 条边界条件结果存在 & 保留 conflict failure 表述。 \\
		听测统计 & 26 participants / 910 ratings 报告存在 & 继续作为主观补充，避免强外推。 \\
		P0 E2/E3 & 报告和脚本存在，原始 outputs 缺失 & 找回或按 runbook 复现。 \\
		Protocol-B 与 legacy report & 已降级为历史材料 & 不进入主论文证据链。 \\
		\bottomrule
	\end{tabular}}
\end{table}

\begin{table}[H]
	\centering
	\caption{4--6 周 NeurIPS main-track 执行节奏。}
	\label{tab:four-six-week-plan}
	\scriptsize
	\resizebox{\textwidth}{!}{%
	\begin{tabular}{p{0.12\textwidth}p{0.22\textwidth}p{0.34\textwidth}p{0.24\textwidth}}
		\toprule
		周次 & KaiMing gate & 预期产物 & 不通过时的处理 \\
		\midrule
		Week 1 & Evidence freeze & canonical manifest、checksums、204 split lock、logger schema v0 & 不写新 claim，先修证据台账。 \\
		Week 2 & Strong baseline parity & PaSST、PANNs、direct regression、Text2FX-style baseline 的同口径结果 & 若任一强基线缺失，主文表标记为 incomplete。 \\
		Week 3 & TRR ablation & Gram/mean/layer/projection/normalization/top-k 消融表 & 若消融不能解释优势，方法贡献降级为 empirical observation。 \\
		Week 4 & Leakage and robustness & near-duplicate audit、真实音频退化 E3、失败类型清单 & 若严格 split 下优势消失，重新定义任务或 claim。 \\
		Week 5 & Metric-perception link & objective--listening correlation、anchor/rater metadata、附录统计 & 若相关性弱，主文把参数距离限定为 controllability proxy。 \\
		Week 6 & Submission freeze & anonymized code/data package、checklist、appendix、paper table lock & 未通过则进入 workshop/technical report 路线。 \\
		\bottomrule
	\end{tabular}}
\end{table}

表~\ref{tab:project-progress} 中最需要优先处理的是 P0、主证据冻结和 logger schema。若目标是正式投稿，建议先把 204 条 audio-grouped Protocol-A、Protocol-C 和 listening test 的 canonical 文件纳入一个干净分支，并在 README 或 manifest 中记录生成命令、数据集快照、统计方法和校验和。只有当 P0 原始 CSV/JSON 被找回或复现后，P0 中的 PaSST/PANNs 与真实退化 E3 结果才适合进入正式论文表格。表~\ref{tab:four-six-week-plan} 给出的节奏属于监督计划，尚未构成已完成 repo fact；每周结束必须用 logger artifact 与统计报告更新本章。

当前报告本身也应视为 single-source working report。它已经根据现有代码和实验产物生成正文，并开始承担 supervisor ledger 的职责，但仍需要人工做三类审阅。第一类是科学审阅，确认是否接受“TRR 是当前已评估 retrieval prior 中最强”这一保守主张。第二类是证据审阅，确认所有主表数值是否应以 audio-grouped 204 口径为唯一来源，并决定是否将 211 口径全部移入附录或历史说明。第三类是投稿审阅，确认 NeurIPS main track 的 baseline、消融、复现、匿名代码和 checklist 是否满足主会标准。
